\documentclass{article}

\usepackage[utf8]{inputenc}     % pdfLaTeX
\usepackage[T1]{fontenc}
\usepackage[magyar]{babel}
\usepackage{amsmath}

% Set page size and margins
% Replace `letterpaper' with `a4paper' for UK/EU standard size
\usepackage[a4paper,top=2cm,bottom=2cm,left=3cm,right=3cm,marginparwidth=1.75cm]{geometry}
%for arrows
\usepackage{textcomp}

\title{Módosítószók}
\author{Lengyel Zoltán Péter\thanks{Csukás Ibolya tanítása alapján}}
\date{Legutóbb frissítve: 2026. Március 20.}

\begin{document}

\maketitle
\tableofcontents
\newpage

\begin{center}
    \section{Módosítószók}
\end{center}
\begin{itemize}
    \item A mondat állításához kommunikációs, érzelmi vagy értékelő jelentéselemeket tesznek hozzá
    
    \item Jelölhetik: 
    \begin{itemize}
        \item A mondat fajtáját: -e, vajon, ugye, bár, bárcsak, de, hadd
    \end{itemize}
    
    \item Árnyalják a mondat jelentését: \begin{itemize}
        \item hiszen, ám, inkább, csak, már
    \end{itemize}
    
    \item Fokozati skálán helyezik el a tulajdonságot: 
    \begin{itemize}
        \item elég, csupa
    \end{itemize}
    
    \item Megbecsülhetnek egy tulajdonságot vagy mennyiséget: 
    \begin{itemize}
        \item körülbelül, alig
    \end{itemize}
    
    \item Kiemelhetik a következő elemet: \begin{itemize}
        \item csak
    \end{itemize}
\end{itemize}

\subsection{Tagadószók és tiltószók}
\begin{itemize}
    \item Tagadják a mondat állításának részét vagy egészét
    \begin{itemize}
        \item Nem kérem,
        \item Nem azt kérem,
        \item Ne vedd el,
        \item Ne azt vedd el.
    \end{itemize}
    
\end{itemize}

\end{document}